% ASSERT_CONVENTION: natural_units=natural, metric_signature=mostly_minus, fourier_convention=physics, coupling_convention=alpha_s, generator_normalization=T(fund)=1/2, covariant_derivative_sign=D_mu=partial_mu-igA
%
% Exercise Set 7: The Landscape of SUSY Dualities
% Part of: The Dualised Standard Model -- Part III Exercises
%
% Conventions (Seiberg hep-th/9411149 + Intriligator-Seiberg):
%   T(fund SU(N)) = 1/2 per Weyl fermion
%   T(vector SO(N)) = 1 (NOT 1/2)
%   T(fund Sp(N)) = 1/2
%   Dual Coxeter numbers: h(SU(N)) = N, h(SO(N)) = N-2, h(Sp(N)) = N+1
%   Sp(N_c) = rank N_c, fundamental dimension 2N_c (Intriligator-Seiberg)
%   A(fund) = 1 (cubic anomaly coefficient per fundamental Weyl fermion)
%   R[Q] determined by anomaly-free condition
%
\documentclass[11pt,a4paper]{article}
% ASSERT_CONVENTION: natural_units=natural, metric_signature=mostly_minus, fourier_convention=physics, coupling_convention=alpha_s, generator_normalization=T(fund)=1/2, covariant_derivative_sign=D_mu=partial_mu-igA
%
% CONVENTIONS (Seiberg hep-th/9411149)
% Metric: (+,-,-,-) (mostly minus)
% T(fund) = 1/2 per Weyl fermion
% b_0 = (11N_c - 2N_f)/3 for N_f Dirac flavours
% D_mu = partial_mu - ig A_mu^a T^a
% A(fund) = 1 (cubic anomaly coefficient per fundamental Weyl fermion)
% alpha_s = g_s^2 / (4 pi)
% R[Q] = (N_f - N_c)/N_f; R[W] = 2
% N_f counts Dirac flavour pairs (each = 2 Weyl fermions)
%
% This file is \input{} by all lecture files.

%% ---------- Document class ----------
% (loaded by the wrapper; preamble only provides packages and macros)

%% ---------- Standard packages ----------
\usepackage{amsmath,amssymb,amsthm,mathtools}
\usepackage{hyperref}
\usepackage[capitalise,noabbrev]{cleveref}
\usepackage{enumitem}
\usepackage{booktabs}
\usepackage{array}
\usepackage{xcolor}
\usepackage{tikz}
\usetikzlibrary{arrows.meta,decorations.markings,positioning,calc}
\usepackage{fancybox}
\usepackage{tcolorbox}
\tcbuselibrary{skins,breakable}
\usepackage{slashed}              % Feynman slash notation
\usepackage{cite}                 % compressed citation lists

%% ---------- Hyperref configuration ----------
\hypersetup{
  colorlinks=true,
  linkcolor=blue!70!black,
  citecolor=green!50!black,
  urlcolor=blue!80!black,
  pdfauthor={Part III Lecture Notes},
  pdftitle={The Dualised Standard Model}
}

%% ---------- Theorem environments ----------
\theoremstyle{plain}
\newtheorem{theorem}{Theorem}[section]
\newtheorem{lemma}[theorem]{Lemma}
\newtheorem{proposition}[theorem]{Proposition}
\newtheorem{corollary}[theorem]{Corollary}

\theoremstyle{definition}
\newtheorem{definition}[theorem]{Definition}
\newtheorem{example}[theorem]{Example}
\newtheorem{exercise}{Exercise}[section]

\theoremstyle{remark}
\newtheorem{remark}[theorem]{Remark}

%% ---------- Physics macros: gauge groups and representations ----------
\newcommand{\Nc}{N_{\!c}}
\newcommand{\Nf}{N_{\!f}}
\newcommand{\SU}[1]{\mathrm{SU}(#1)}
\newcommand{\UU}[1]{\mathrm{U}(#1)}

% Representations
\newcommand{\fund}{\square}
\newcommand{\antifund}{\overline{\square}}
\newcommand{\adj}{\mathrm{adj}}
\newcommand{\antisym}{\wedge^{\!2}}        % antisymmetric rank-2
\newcommand{\sym}{\mathrm{S}_2}           % symmetric rank-2

%% ---------- Physics macros: group theory constants ----------
% Dynkin index: Tr(T^a T^b) = T(R) delta^{ab}
\newcommand{\Tfund}{\tfrac{1}{2}}          % T(fund) = 1/2 per Weyl fermion
\newcommand{\Tadj}{\Nc}                    % T(adj) = N_c
\newcommand{\Cfund}{\frac{\Nc^2 - 1}{2\Nc}}  % C_2(fund)
\newcommand{\Cadj}{\Nc}                    % C_2(adj)

% Cubic anomaly coefficient: A(R) = Tr_R[T^a {T^b, T^c}] with A(fund) = 1
\newcommand{\Afund}{1}

%% ---------- Physics macros: beta function ----------
\newcommand{\bzero}{b_0}
\newcommand{\bone}{b_1}
\newcommand{\alphas}{\alpha_s}
\newcommand{\gs}{g_s}

%% ---------- Physics macros: SUSY / R-charge ----------
\newcommand{\Rcharge}[1]{R[#1]}
\newcommand{\Wsup}{W}                      % superpotential

%% ---------- Physics macros: covariant derivative and field strength ----------
\newcommand{\Dmu}{D_\mu}
\newcommand{\Fmn}{F_{\mu\nu}}
\newcommand{\Ftilde}{\widetilde{F}_{\mu\nu}}

%% ---------- Physics macros: common abbreviations ----------
\newcommand{\ie}{\textit{i.e.}}
\newcommand{\eg}{\textit{e.g.}}
\newcommand{\cf}{\textit{cf.}}
\newcommand{\IR}{\ensuremath{\mathrm{IR}}}
\newcommand{\UV}{\ensuremath{\mathrm{UV}}}
\newcommand{\AF}{\ensuremath{\mathrm{AF}}}
\newcommand{\BZ}{\ensuremath{\mathrm{BZ}}}
\newcommand{\NSVZ}{\ensuremath{\mathrm{NSVZ}}}
\newcommand{\SQCD}{\ensuremath{\mathrm{SQCD}}}
\newcommand{\QCD}{\ensuremath{\mathrm{QCD}}}
\newcommand{\SM}{\ensuremath{\mathrm{SM}}}
\newcommand{\Tr}{\mathrm{Tr}}

%% ---------- Convention box macro ----------
\newtcolorbox{conventionbox}[1][]{%
  enhanced,
  colback=blue!5!white,
  colframe=blue!60!black,
  fonttitle=\bfseries,
  title={Conventions},
  #1
}

%% ---------- Comparison table style ----------
\newtcolorbox{comparisontable}[1][]{%
  enhanced,
  colback=orange!5!white,
  colframe=orange!70!black,
  fonttitle=\bfseries,
  title={Comparison},
  #1
}

%% ---------- Warning / important box ----------
\newtcolorbox{warningbox}[1][]{%
  enhanced,
  colback=red!5!white,
  colframe=red!70!black,
  fonttitle=\bfseries,
  title={Important},
  #1
}

%% ---------- Preview / forward-reference box ----------
\newtcolorbox{previewbox}[1][]{%
  enhanced,
  colback=green!5!white,
  colframe=green!50!black,
  fonttitle=\bfseries,
  title={Preview},
  #1
}

%% ---------- Page layout ----------
\usepackage[margin=2.5cm]{geometry}
\setlength{\parskip}{0.4em}
\setlength{\parindent}{0em}

%% ---------- Section formatting ----------
\usepackage{titlesec}
\titleformat{\section}{\Large\bfseries}{\thesection}{1em}{}
\titleformat{\subsection}{\large\bfseries}{\thesubsection}{1em}{}

%% ---------- End of preamble ----------


%% ---------- Additional macros for this exercise ----------
\newcommand{\SO}[1]{\mathrm{SO}(#1)}
\newcommand{\Sp}[1]{\mathrm{Sp}(#1)}
\newcommand{\Nctilde}{\widetilde{N}_{\!c}}

\title{\textbf{Exercise Set 7: The Landscape of SUSY Dualities}}
\author{The Dualised Standard Model --- Part III Exercises}
\date{Lent Term 2027}

\begin{document}
\maketitle

\begin{conventionbox}[title={Conventions for these exercises}]
\renewcommand{\arraystretch}{1.3}
\begin{tabular}{@{}l l@{}}
\toprule
\textbf{Quantity} & \textbf{Convention} \\
\midrule
Units & $\hbar = c = 1$ (natural) \\
Metric & $\eta_{\mu\nu} = \mathrm{diag}(+1,-1,-1,-1)$ (mostly minus) \\
$T(\text{fund } \SU{N})$ & $\tfrac{1}{2}$ per Weyl fermion \\
$T(\text{vector } \SO{N})$ & $1$ (Intriligator--Seiberg convention) \\
$T(\text{fund } \Sp{N})$ & $\tfrac{1}{2}$ \\
$\Sp{\Nc}$ & rank $\Nc$; fund.\ dim.\ $2\Nc$; \; $\Sp{1} = \SU{2}$ \\
Dual Coxeter & $h_{\SU{N}} = N$, \; $h_{\SO{N}} = N-2$, \; $h_{\Sp{N}} = N+1$ \\
$A(\fund)$ & $1$ (cubic anomaly coefficient per fundamental Weyl fermion) \\
\bottomrule
\end{tabular}

\medskip

\noindent References: Lecture~9 (all sections);
Intriligator--Seiberg~\cite{IS9506101};
Kutasov--Schwimmer~\cite{KS9505004};
Cs\'aki--Schmaltz--Skiba~\cite{CSS9612207}.

\medskip

\noindent\textbf{Total marks: 65. \quad Suggested time: 90 minutes.}
\end{conventionbox}

\bigskip


%% ========================================================================
%%  PROBLEM 1: SO(N_c) Anomaly Matching
%% ========================================================================
\section*{Problem 1: Anomaly Matching in $\SO{\Nc}$ Duality
  \hfill $\star\star\star$ \normalsize(40 min, 30 marks)}
\addcontentsline{toc}{section}{Problem 1: $\SO{\Nc}$ anomaly matching}
\label{prob:SO-anomaly}

\noindent\textbf{Background.}
Intriligator and Seiberg~\cite{IS9506101} established a duality for
$\mathcal{N}=1$ $\SO{\Nc}$ gauge theory with $\Nf$ chiral superfields
$Q_i$ ($i = 1, \ldots, \Nf$) in the vector representation (dimension
$\Nc$) and no tree-level superpotential.  The dual (magnetic) theory
has gauge group $\SO{\Nf - \Nc + 4}$, with $\Nf$ dual vectors $q_i$
and a gauge-singlet symmetric meson matrix $M_{ij} = M_{ji}$
(mapping to $Q_i Q_j$ in the electric theory), with magnetic
superpotential $W_{\text{mag}} = M_{ij}\,q_i\,q_j / (2\mu)$.

The conformal window is $\tfrac{3}{2}(\Nc - 2) < \Nf < 3(\Nc - 2)$.

\medskip

\noindent\textbf{Questions.}
\begin{enumerate}[label=(\alph*)]

  \item \textbf{[5 marks]}
    Consider $\SO{10}$ with $\Nf = 13$ vectors.
    \begin{enumerate}[label=(\roman*)]
      \item Verify that $\Nf = 13$ lies inside the conformal window for
        $\SO{10}$.
      \item State the dual gauge group and its vector representation
        dimension.
      \item What is the dimension of the symmetric meson $M_{ij}$ as a
        matrix?  How many independent components does it have?
    \end{enumerate}

  \item \textbf{[10 marks]}
    Compute the mixed anomaly coefficient
    $\mathcal{A}[\SU{\Nf}^2\,\UU{1}_R]$ on the \textbf{electric} side.

    \begin{enumerate}[label=(\roman*)]
      \item The non-anomalous $R$-charge of the electric quark is
        determined by $T(\adj) + \Nf\,T(\text{vector})\,(R[Q] - 1) = 0$.
        Use $T(\adj_{\SO{\Nc}}) = h_{\SO{\Nc}} = \Nc - 2$ and
        $T(\text{vector}) = 1$ to find $R[Q]$ for general $\Nc$, $\Nf$.

      \item The fermionic component's $R$-charge is
        $R_{\text{ferm}}[Q] = R[Q] - 1$.  The mixed anomaly receives
        a contribution:
        \[
          \mathcal{A}^{\text{el}}
          = \dim(\text{gauge rep of } Q)\;\cdot\;
            T(\fund_{\SU{\Nf}})\;\cdot\; R_{\text{ferm}}[Q]\,.
        \]
        Evaluate this for general $\Nc$, $\Nf$, and then for
        $(\Nc, \Nf) = (10, 13)$.
    \end{enumerate}

  \item \textbf{[10 marks]}
    Compute $\mathcal{A}[\SU{\Nf}^2\,\UU{1}_R]$ on the
    \textbf{magnetic} side.
    \begin{enumerate}[label=(\roman*)]
      \item Write down the dual Coxeter number $h_{\text{mag}}$ of the
        magnetic gauge group.  Hence find $R[q]$ and
        $R_{\text{ferm}}[q]$.
      \item From $M_{ij} \sim Q_i Q_j$, determine $R[M]$ and
        $R_{\text{ferm}}[M]$.
      \item The meson is a symmetric matrix under $\SU{\Nf}$,
        transforming in the $\mathrm{S}_2$ representation with Dynkin
        index $T(\mathrm{S}_2\;\SU{\Nf}) = (\Nf + 2)/2$.  Compute the
        two contributions to the magnetic anomaly:
        \[
          \mathcal{A}^{\text{mag}}_q
          = \dim(\text{magnetic vector})\;\cdot\;
            T(\fund_{\SU{\Nf}})\;\cdot\; R_{\text{ferm}}[q]\,,
        \]
        \[
          \mathcal{A}^{\text{mag}}_M
          = 1 \;\cdot\;
            T(\mathrm{S}_2\;\SU{\Nf})\;\cdot\; R_{\text{ferm}}[M]\,.
        \]
        Sum them.
    \end{enumerate}

  \item \textbf{[5 marks]}
    Verify that $\mathcal{A}^{\text{el}} = \mathcal{A}^{\text{mag}}$
    both in general (for arbitrary $\Nc$, $\Nf$) and numerically for
    $(\Nc, \Nf) = (10, 13)$.

\end{enumerate}

\medskip

\noindent\textit{Hint:} The algebra in part~(c)(iii) involves expanding
$(\Nf - \Nc + 4)(\Nf - \Nc + 2)$ and $(\Nf + 2)(\Nf - 2\Nc + 4)$.
Keep both terms over a common denominator $2\Nf$ and simplify.


\bigskip
\hrule
\bigskip


%% ========================================================================
%%  SOLUTION TO PROBLEM 1
%% ========================================================================
\begin{proof}[\textbf{Solution to Problem 1}]

\textbf{(a) Setup for $\SO{10}$, $\Nf = 13$.}

\textit{(i) Conformal window check.}

For $\SO{10}$: $h = \Nc - 2 = 8$.  The conformal window is:
\[
  \frac{3 \times 8}{2} < \Nf < 3 \times 8
  \qquad\Longrightarrow\qquad
  12 < \Nf < 24\,.
\]
Since $12 < 13 < 24$, we have $\Nf = 13$ inside the conformal window.
\quad$\checkmark$

\textit{(ii) Dual gauge group.}

The magnetic gauge group is:
\[
  \SO{\Nf - \Nc + 4} = \SO{13 - 10 + 4} = \SO{7}\,.
\]
The vector representation of $\SO{7}$ has dimension $7$.

\textit{(iii) Meson dimension.}

$M_{ij}$ is a $13 \times 13$ symmetric matrix.  The number of
independent components is $\Nf(\Nf + 1)/2 = 13 \times 14/2 = 91$.

\bigskip

\textbf{(b) Electric-side anomaly.}

\textit{(i) $R$-charge.}

From $T(\adj) + \Nf\,T(\text{vector})\,(R[Q]-1) = 0$:
\[
  (\Nc - 2) + \Nf \cdot 1 \cdot (R[Q] - 1) = 0\,,
\]
\[
  R[Q] = 1 - \frac{\Nc - 2}{\Nf}\,.
\]
For $(\Nc, \Nf) = (10, 13)$:
$R[Q] = 1 - 8/13 = 5/13$.

\textit{(ii) Mixed anomaly.}

The fermionic $R$-charge:
\[
  R_{\text{ferm}}[Q] = R[Q] - 1 = -\frac{\Nc - 2}{\Nf}\,.
\]
For $(10, 13)$: $R_{\text{ferm}}[Q] = -8/13$.

The electric quarks $Q_i$ transform in the vector of $\SO{\Nc}$
(dimension $\Nc$) and the fundamental of $\SU{\Nf}$ (with Dynkin
index $T(\fund) = 1/2$):
\begin{align}
  \mathcal{A}^{\text{el}}
  &= \Nc \cdot \frac{1}{2} \cdot
    \biggl(-\frac{\Nc - 2}{\Nf}\biggr)
  = -\frac{\Nc(\Nc - 2)}{2\Nf}\,. \label{eq:Ael-SO}
\end{align}

For $(\Nc, \Nf) = (10, 13)$:
\[
  \mathcal{A}^{\text{el}}
  = -\frac{10 \times 8}{2 \times 13}
  = -\frac{80}{26}
  = -\frac{40}{13}\,.
\]
\[
  \boxed{\mathcal{A}^{\text{el}} = -\frac{40}{13}\,.}
\]

\bigskip

\textbf{(c) Magnetic-side anomaly.}

\textit{(i) Dual quarks.}

The magnetic gauge group is $\SO{\Nf - \Nc + 4}$ with dual Coxeter
number:
\[
  h_{\text{mag}} = (\Nf - \Nc + 4) - 2 = \Nf - \Nc + 2\,.
\]
For $(10, 13)$: $h_{\text{mag}} = 13 - 10 + 2 = 5$.

The $R$-charge of the magnetic quark:
\[
  R[q] = 1 - \frac{h_{\text{mag}}}{\Nf}
  = 1 - \frac{\Nf - \Nc + 2}{\Nf}
  = \frac{\Nc - 2}{\Nf}\,.
\]
For $(10,13)$: $R[q] = 8/13$.

The fermionic $R$-charge:
\[
  R_{\text{ferm}}[q] = R[q] - 1 = \frac{\Nc - 2 - \Nf}{\Nf}\,.
\]
For $(10,13)$: $R_{\text{ferm}}[q] = (8 - 13)/13 = -5/13$.

\textit{(ii) Meson.}

From $M_{ij} \sim Q_i Q_j$:
\[
  R[M] = 2\,R[Q] = \frac{2(\Nf - \Nc + 2)}{\Nf}\,.
\]
For $(10,13)$: $R[M] = 2 \times 5/13 = 10/13$.

\[
  R_{\text{ferm}}[M] = R[M] - 1 = \frac{\Nf - 2\Nc + 4}{\Nf}
  = \frac{\Nf - 2(\Nc - 2)}{\Nf}\,.
\]
For $(10,13)$: $R_{\text{ferm}}[M] = (13 - 16)/13 = -3/13$.

\textit{(iii) Magnetic anomaly contributions.}

The dual quarks are in the vector of $\SO{7}$ (dimension $7$) and
the fundamental of $\SU{13}$:
\begin{align}
  \mathcal{A}^{\text{mag}}_q
  &= (\Nf - \Nc + 4) \cdot \frac{1}{2}
    \cdot \frac{\Nc - 2 - \Nf}{\Nf}\,.
\end{align}

For $(10,13)$:
\[
  \mathcal{A}^{\text{mag}}_q
  = 7 \cdot \frac{1}{2} \cdot \frac{-5}{13}
  = -\frac{35}{26}\,.
\]

The meson is a gauge singlet (dimension 1 under the magnetic gauge
group), in the symmetric $\mathrm{S}_2$ representation of $\SU{13}$
with $T(\mathrm{S}_2) = (13+2)/2 = 15/2$:
\begin{align}
  \mathcal{A}^{\text{mag}}_M
  &= 1 \cdot \frac{\Nf + 2}{2}
    \cdot \frac{\Nf - 2\Nc + 4}{\Nf}\,.
\end{align}

For $(10,13)$:
\[
  \mathcal{A}^{\text{mag}}_M
  = 1 \cdot \frac{15}{2} \cdot \frac{-3}{13}
  = -\frac{45}{26}\,.
\]

\textbf{Total magnetic anomaly:}
\[
  \mathcal{A}^{\text{mag}}
  = -\frac{35}{26} - \frac{45}{26}
  = -\frac{80}{26}
  = -\frac{40}{13}\,.
\]
\[
  \boxed{\mathcal{A}^{\text{mag}} = -\frac{40}{13}\,.}
\]

\bigskip

\textbf{(d) Verification of anomaly matching.}

\textit{Numerical check:}
$\mathcal{A}^{\text{el}} = -40/13 = \mathcal{A}^{\text{mag}}$.
\quad$\checkmark$

\textit{General proof.}  We need to show:
\[
  \mathcal{A}^{\text{mag}}
  = \frac{1}{2\Nf}\Bigl[
    -(\Nf - \Nc + 4)(\Nf - \Nc + 2)
    + (\Nf + 2)(\Nf - 2\Nc + 4)
  \Bigr]
  = -\frac{\Nc(\Nc - 2)}{2\Nf}\,.
\]

Expand the first product:
\begin{align*}
  (\Nf - \Nc + 4)(\Nf - \Nc + 2)
  &= \Nf^2 - 2\Nc\Nf + \Nc^2 + 6\Nf - 6\Nc + 8\,.
\end{align*}

Expand the second product:
\begin{align*}
  (\Nf + 2)(\Nf - 2\Nc + 4)
  &= \Nf^2 - 2\Nc\Nf + 4\Nf + 2\Nf - 4\Nc + 8 \\
  &= \Nf^2 - 2\Nc\Nf + 6\Nf - 4\Nc + 8\,.
\end{align*}

Subtracting:
\begin{align*}
  &-(\Nf^2 - 2\Nc\Nf + \Nc^2 + 6\Nf - 6\Nc + 8) \\
  &\quad + (\Nf^2 - 2\Nc\Nf + 6\Nf - 4\Nc + 8) \\
  &= -\Nc^2 + 6\Nc - 4\Nc \\
  &= -\Nc^2 + 2\Nc = -\Nc(\Nc - 2)\,.
\end{align*}

Therefore:
\[
  \mathcal{A}^{\text{mag}}
  = \frac{-\Nc(\Nc - 2)}{2\Nf}
  = \mathcal{A}^{\text{el}}\,.
  \quad\checkmark
\]

The anomaly matching holds for all $\Nc$ and $\Nf$, confirming the
Intriligator--Seiberg duality~\cite{IS9506101}.

\end{proof}


\bigskip
\hrule
\bigskip


%% ========================================================================
%%  PROBLEM 2: Kutasov-Schwimmer Duality
%% ========================================================================
\section*{Problem 2: Kutasov--Schwimmer Duality and the $k = 1$ Limit
  \hfill $\star\star$ \normalsize(25 min, 20 marks)}
\addcontentsline{toc}{section}{Problem 2: Kutasov--Schwimmer duality}
\label{prob:KS-duality}

\noindent\textbf{Background.}
Kutasov and Schwimmer~\cite{KS9505004} generalised Seiberg duality to
$\SU{\Nc}$ with $\Nf$ fundamental flavours $(Q_i, \widetilde{Q}_i)$ and
one adjoint chiral superfield $\Phi$, with tree-level superpotential
$W_{\text{el}} = g\,\Tr(\Phi^{k+1})/(k+1)$ for a positive integer $k$.

The magnetic theory has gauge group $\SU{k\Nf - \Nc}$, dual quarks
$(q_i, \widetilde{q}^{\,i})$, a dual adjoint $\phi$, and $k$ meson
species $M_j = Q\,\Phi^j\,\widetilde{Q}$ for $j = 0, 1, \ldots, k-1$.

The conformal window is $3\Nc/(k+1) < \Nf < 3\Nc/k$.

\medskip

\noindent\textbf{Questions.}
\begin{enumerate}[label=(\alph*)]

  \item \textbf{[5 marks]}
    Derive the conformal window for the Kutasov--Schwimmer theory.
    \begin{enumerate}[label=(\roman*)]
      \item The one-loop beta function coefficient for $\SU{\Nc}$ with
        $\Nf$ fundamentals and one adjoint is
        $b_0 = 3\Nc - \Nf - \Nc = 2\Nc - \Nf$.  Verify this
        using $T(\adj) = \Nc$ and $T(\fund) = 1/2$.
      \item At the fixed point, the superpotential
        $W = \Tr(\Phi^{k+1})$ imposes $R[\Phi] = 2/(k+1)$.  Using
        the anomaly-free condition (with $\Phi$ carrying $R$-charge
        $2/(k+1)$ rather than the free value), the upper bound on $\Nf$
        from asymptotic freedom of the electric theory is
        $\Nf < 3\Nc/k$.  Show this.
      \item The lower bound comes from requiring the magnetic theory
        $\SU{k\Nf - \Nc}$ to also be asymptotically free.  Show that
        this gives $\Nf > 3\Nc/(k+1)$.
    \end{enumerate}

  \item \textbf{[5 marks]}
    Verify that the conformal window formula reduces to the standard
    Seiberg result $\tfrac{3}{2}\Nc < \Nf < 3\Nc$ when $k = 1$.

  \item \textbf{[5 marks]}
    Show that the $k = 1$ case of the KS duality reduces to Seiberg
    duality by following these steps:
    \begin{enumerate}[label=(\roman*)]
      \item At $k = 1$, the electric superpotential is
        $W_{\text{el}} = (g/2)\,\Tr(\Phi^2)$.  What is the physical
        interpretation of this term?  What happens to $\Phi$ in the IR?
      \item Show that the magnetic gauge group at $k = 1$ is
        $\SU{\Nf - \Nc}$, matching Lecture~3.
      \item Show that there is exactly one meson species $M_0 =
        Q\,\widetilde{Q}$ and that the magnetic superpotential reduces
        to $W_{\text{mag}} = M\,q\,\widetilde{q}/\mu$ after
        integrating out the massive dual adjoint $\phi$.
    \end{enumerate}

  \item \textbf{[5 marks]}
    Consider $\SU{4}$ with $k = 2$ and $\Nf = 5$.
    \begin{enumerate}[label=(\roman*)]
      \item Verify that $\Nf = 5$ lies inside the conformal window.
      \item State the magnetic gauge group.
      \item How many meson species $M_j$ are there?  What are their
        electric-theory definitions?
    \end{enumerate}

\end{enumerate}

\bigskip
\hrule
\bigskip


%% ========================================================================
%%  SOLUTION TO PROBLEM 2
%% ========================================================================
\begin{proof}[\textbf{Solution to Problem 2}]

\textbf{(a) Deriving the conformal window.}

\textit{(i) Beta function coefficient.}

For $\SU{\Nc}$ with $\Nf$ fundamentals (each with $T(\fund) = 1/2$),
$\Nf$ antifundamentals (same), and one adjoint ($T(\adj) = \Nc$):
\begin{align*}
  b_0 &= \frac{11}{3}\,T(\adj) - \frac{2}{3}\,\bigl(2\Nf\,T(\fund)
    + T(\adj)\bigr) \\
  &= \frac{11}{3}\,\Nc - \frac{2}{3}\bigl(\Nf + \Nc\bigr)
  = 3\Nc - \Nf - \Nc \quad\text{(after simplification)}\\
  &= 2\Nc - \Nf\,.
  \quad\checkmark
\end{align*}

Here we counted $2\Nf$ Weyl fermions from the $\Nf$ pairs
$(Q_i, \widetilde{Q}_i)$ and one Weyl adjoint from $\Phi$.

\textit{(ii) Upper bound.}

At the fixed point, $R[\Phi] = 2/(k+1)$ is fixed by the superpotential.
The anomalous dimension of $\Phi$ is $\gamma_\Phi = R[\Phi] - 2/3 \cdot 1$
(in the standard relation $\dim(\Phi) = 1 + \gamma/2$, with
$R[\Phi] = 2/3(1+\gamma_\Phi/2)$).

More directly: the loss of asymptotic freedom occurs when the effective
$b_0$ (incorporating the fixed-point anomalous dimensions) vanishes.
The NSVZ beta function at the fixed point requires:
\[
  3\Nc - \Nf\,(1 - \gamma_Q) - \Nc\,(1 - \gamma_\Phi) = 0\,.
\]
The superpotential $W = \Tr(\Phi^{k+1})$ has $R$-charge~$2$, so
$(k+1)\,R[\Phi] = 2$, giving $R[\Phi] = 2/(k+1)$.  The marginal
coupling condition for the superpotential vertex requires that
$\Phi$ has dimension $1$ at the fixed point when $k + 1 = 2$ (i.e.\
$k=1$), and in general $\gamma_\Phi = 2R[\Phi]/2 - 2/3$.

The key constraint is that the $\SU{\Nf}^3$ anomaly matching combined
with the fixed-point unitarity bound requires~\cite{KS9505004,KSS9507040}:
\[
  \Nf < \frac{3\Nc}{k}\,.
\]
To see this simply: the superpotential constraint fixes
$\gamma_\Phi^* = -2/(k+1)$ (so that $R[\Phi]$ at the fixed point is
$2/(k+1)$, and $\Delta(\Phi) = 3R[\Phi]/2 = 3/(k+1)$).  The
$a$-maximisation condition (or equivalently the
unitarity bound $R[Q] \geq 1/3$ at the fixed point for gauge-invariant
operators) then yields $\Nf < 3\Nc/k$.
\quad$\checkmark$

\textit{(iii) Lower bound.}

The magnetic theory has gauge group $\SU{k\Nf - \Nc}$, with $\Nf$
pairs of fundamental/antifundamental quarks and one adjoint $\phi$.
Its beta function coefficient is:
\[
  b_0^{\text{mag}} = 2(k\Nf - \Nc) - \Nf\,.
\]
Asymptotic freedom requires $b_0^{\text{mag}} > 0$, which gives
$\Nf > 3\Nc/(k+1)$.

More precisely, the same fixed-point constraint applied to the magnetic
theory (with the same superpotential structure) yields:
\[
  \Nf > \frac{3\Nc}{k+1}\,,
\]
which is the lower bound.  Combining:
\[
  \boxed{\frac{3\Nc}{k+1} < \Nf < \frac{3\Nc}{k}\,.}
  \quad\checkmark
\]

\bigskip

\textbf{(b) The $k = 1$ limit of the conformal window.}

Setting $k = 1$:
\[
  \frac{3\Nc}{k+1} < \Nf < \frac{3\Nc}{k}
  \quad\Longrightarrow\quad
  \frac{3\Nc}{2} < \Nf < 3\Nc\,.
\]
This is exactly the Seiberg conformal window from Lecture~3.
\quad$\checkmark$

\bigskip

\textbf{(c) The $k = 1$ case reduces to Seiberg duality.}

\textit{(i) Electric superpotential at $k = 1$.}

$W_{\text{el}} = (g/2)\,\Tr(\Phi^2)$.  This is a \textbf{mass term}
for the adjoint field: $m_\Phi = g$.  Below the energy scale
$\mu \ll g$, the adjoint $\Phi$ is integrated out, leaving $\SU{\Nc}$
SQCD with $\Nf$ flavours and no superpotential---exactly the electric
theory of Lecture~3.

\textit{(ii) Magnetic gauge group at $k = 1$.}

\[
  \SU{k\Nf - \Nc}\big|_{k=1} = \SU{\Nf - \Nc}\,.
  \quad\checkmark
\]

\textit{(iii) Meson content and superpotential at $k = 1$.}

With $k = 1$, we have one meson ($j = 0$ only):
\[
  M_0 = Q\,\Phi^0\,\widetilde{Q} = Q\,\widetilde{Q} = M\,,
\]
which is the standard Seiberg meson.

The dual adjoint $\phi$ also receives a mass term
$W \supset (g'/2)\,\Tr(\phi^2)$.  Integrating it out, the magnetic
superpotential becomes:
\[
  W_{\text{mag}}
  \;\xrightarrow{\;\text{integrate out}\;\phi\;}
  \;\frac{1}{\mu}\, M\, q\, \widetilde{q}\,.
\]
This is the Seiberg magnetic superpotential from Lecture~3.
\quad$\checkmark$

\bigskip

\textbf{(d) Explicit example: $\SU{4}$, $k = 2$, $\Nf = 5$.}

\textit{(i) Conformal window.}

\[
  \frac{3 \times 4}{3} < \Nf < \frac{3 \times 4}{2}
  \quad\Longrightarrow\quad
  4 < \Nf < 6\,.
\]
Since $4 < 5 < 6$, $\Nf = 5$ is inside the conformal window.
\quad$\checkmark$

\textit{(ii) Magnetic gauge group.}

\[
  \SU{k\Nf - \Nc} = \SU{2 \times 5 - 4} = \SU{6}\,.
\]

\textit{(iii) Meson species.}

With $k = 2$, there are $k = 2$ meson species ($j = 0, 1$):
\begin{align*}
  M_0 &= Q\,\Phi^0\,\widetilde{Q} = Q\,\widetilde{Q}\,,\\
  M_1 &= Q\,\Phi^1\,\widetilde{Q} = Q\,\Phi\,\widetilde{Q}\,.
\end{align*}
Each is a $5 \times 5$ matrix transforming as
$(\fund, \antifund)$ under $\SU{5}_L \times \SU{5}_R$.

\end{proof}


\bigskip
\hrule
\bigskip


%% ========================================================================
%%  PROBLEM 3: s-Confinement
%% ========================================================================
\section*{Problem 3: s-Confinement as the Boundary of Duality
  \hfill $\star\star$ \normalsize(25 min, 15 marks)}
\addcontentsline{toc}{section}{Problem 3: s-Confinement}
\label{prob:s-confinement}

\noindent\textbf{Background.}
An $\mathcal{N}=1$ gauge theory is \textbf{s-confining} (smoothly
confining) if the infrared degrees of freedom are gauge-invariant
composites with a smooth confining superpotential and no residual gauge
symmetry~\cite{CSS9612207}.  These theories sit at the boundary of the
duality landscape, where the magnetic gauge group has rank zero.

\medskip

\noindent\textbf{Questions.}
\begin{enumerate}[label=(\alph*)]

  \item \textbf{[5 marks]}
    Consider $\SU{\Nc}$ SQCD with $\Nf = \Nc + 1$ flavours.
    \begin{enumerate}[label=(\roman*)]
      \item Show that the Seiberg dual gauge group is
        $\SU{\Nf - \Nc} = \SU{1}$, which is trivial.
      \item The confined degrees of freedom are the meson
        $M^i{}_j = Q^i\,\widetilde{Q}_j$ (an $(\Nc+1) \times (\Nc+1)$
        matrix), the baryon $B_{i_1 \cdots i_{\Nc}}
        = \epsilon^{a_1 \cdots a_{\Nc}} Q^{i_1}_{a_1} \cdots
        Q^{i_{\Nc}}_{a_{\Nc}}$, and the antibaryon
        $\widetilde{B}^{i_1 \cdots i_{\Nc}}$ (defined analogously).
        How many independent degrees of freedom are there?
        (\textit{Hint:} $B$ has $\binom{\Nc+1}{\Nc} = \Nc + 1$
        components.)
    \end{enumerate}

  \item \textbf{[5 marks]}
    The confining superpotential is
    (Seiberg~\cite{Seiberg1995}):
    \begin{equation}\label{eq:s-confine-W}
      W_{\text{conf}} = \frac{1}{\Lambda^{2\Nc - 1}}
      \left(B_i\, M^i{}_j\, \widetilde{B}^j - \det M\right)\,.
    \end{equation}
    \begin{enumerate}[label=(\roman*)]
      \item Verify that $W_{\text{conf}}$ has the correct mass
        dimension: recall that in the SUSY conventions used here,
        the superpotential has mass dimension $3$ (since the
        superspace integral $\int d^2\theta$ has dimension $1$).
        The fields $M$, $B$, $\widetilde{B}$ are dimensionless
        composites in canonical normalisation.  What is the
        dimension of $\Lambda^{2\Nc - 1}$?
      \item Verify that $W_{\text{conf}}$ has $R$-charge $2$.
        Use $R[Q] = 1/(\Nc + 1)$ (from
        $R[Q] = (\Nf - \Nc)/\Nf = 1/(\Nc+1)$).  Compute
        $R[M]$, $R[B]$, $R[\widetilde{B}]$, and the
        $R$-charge of each term in the superpotential.
    \end{enumerate}

  \item \textbf{[5 marks]}
    Explain why s-confinement sits at the \textbf{boundary of the
    duality landscape} by completing this table:

    \renewcommand{\arraystretch}{1.3}
    \begin{center}
    \begin{tabular}{@{}l l l@{}}
    \toprule
    \textbf{Gauge group} & \textbf{s-confining at} &
      \textbf{Dual $\to$ trivial} \\
    \midrule
    $\SU{\Nc}$ & $\Nf = \rule{2cm}{0.4pt}$ &
      $\SU{\rule{2cm}{0.4pt}} = \SU{1}$ \\
    $\Sp{\Nc}$ & $\Nf = \rule{2cm}{0.4pt}$ &
      $\Sp{\rule{2cm}{0.4pt}} = \Sp{0}$ \\
    $\SO{\Nc}$ & $\Nf = \rule{2cm}{0.4pt}$ &
      $\SO{\rule{2cm}{0.4pt}} = \SO{0}$ \\
    \bottomrule
    \end{tabular}
    \end{center}

    For the $\Sp{\Nc}$ case, verify that $\Sp{1}$ with
    $2\Nf = 8$ fundamentals (i.e.\ $\Nf = 4$) is
    \emph{not} the s-confining point but rather the
    \emph{self-dual} point.  What is the s-confining
    number of flavours for $\Sp{1}$?

\end{enumerate}

\bigskip
\hrule
\bigskip


%% ========================================================================
%%  SOLUTION TO PROBLEM 3
%% ========================================================================
\begin{proof}[\textbf{Solution to Problem 3}]

\textbf{(a) $\SU{\Nc}$ with $\Nf = \Nc + 1$.}

\textit{(i) Trivial dual.}

The Seiberg dual gauge group is:
\[
  \SU{\Nf - \Nc} = \SU{(\Nc + 1) - \Nc} = \SU{1}\,.
\]
$\SU{1}$ is the trivial group (no gauge dynamics).  There is no
magnetic gauge symmetry: the theory \textbf{s-confines}.
\quad$\checkmark$

\textit{(ii) Counting degrees of freedom.}

\begin{itemize}
  \item Meson $M^i{}_j$: an $(\Nc + 1) \times (\Nc + 1)$ complex
    matrix, so $(\Nc + 1)^2$ complex degrees of freedom.
  \item Baryon $B_{i_1 \cdots i_{\Nc}}$: totally antisymmetric in
    the $\Nc$ flavour indices chosen from $\{1, \ldots, \Nc + 1\}$.
    The number of independent components is
    $\binom{\Nc + 1}{\Nc} = \Nc + 1$.
  \item Antibaryon $\widetilde{B}^{i_1 \cdots i_{\Nc}}$: same
    counting, $\Nc + 1$ independent components.
\end{itemize}

Total: $(\Nc + 1)^2 + 2(\Nc + 1) = (\Nc + 1)(\Nc + 3)$.

\textbf{Check against the microscopic theory:}
The electric theory has $\Nf \times 2 \times \Nc = 2\Nc(\Nc + 1)$
chiral superfield degrees of freedom (before gauge fixing), minus
$\Nc^2 - 1$ for the gauge degrees of freedom.  The number of
physical moduli should match: $2\Nc(\Nc + 1) - (\Nc^2 - 1)
= \Nc^2 + 2\Nc + 1 = (\Nc + 1)^2$.  With the constraints from the
confining superpotential (which imposes $\Nc + 1$ relations), we
get $(\Nc + 1)^2 - 1$ physical moduli, consistent with the
constrained system of $M$, $B$, $\widetilde{B}$.

\bigskip

\textbf{(b) Dimensional analysis and $R$-charge.}

\textit{(i) Mass dimension.}

In the conventions where $[Q] = 1$ (canonical dimension of a chiral
superfield), the composites have dimensions:
\begin{align*}
  [M^i{}_j] &= [Q\,\widetilde{Q}] = 2\,, \\
  [B] &= [Q^{\Nc}] = \Nc\,, \\
  [\widetilde{B}] &= [\widetilde{Q}^{\Nc}] = \Nc\,.
\end{align*}

The first term $B_i M^i{}_j \widetilde{B}^j$ has dimension
$\Nc + 2 + \Nc = 2\Nc + 2$.  The second term $\det M$ involves
$\Nc + 1$ factors of $M$, so $[\det M] = 2(\Nc + 1) = 2\Nc + 2$.

Both terms have dimension $2\Nc + 2$.  The superpotential must have
dimension $3$, so:
\[
  [W_{\text{conf}}] = 3
  \quad\Longrightarrow\quad
  [\Lambda^{2\Nc - 1}] = (2\Nc + 2) - 3 = 2\Nc - 1\,.
\]
\[
  \boxed{[\Lambda^{2\Nc - 1}] = 2\Nc - 1\,.}
  \quad\checkmark
\]

This is consistent with $[\Lambda] = 1$ (energy).

\textit{(ii) $R$-charge verification.}

With $\Nf = \Nc + 1$:
\[
  R[Q] = \frac{\Nf - \Nc}{\Nf} = \frac{1}{\Nc + 1}\,.
\]

The composites:
\begin{align*}
  R[M] &= R[Q] + R[\widetilde{Q}] = \frac{2}{\Nc + 1}\,, \\
  R[B] &= \Nc\,R[Q] = \frac{\Nc}{\Nc + 1}\,, \\
  R[\widetilde{B}] &= \Nc\,R[\widetilde{Q}] = \frac{\Nc}{\Nc + 1}\,.
\end{align*}

The dynamical scale $\Lambda$ carries $R$-charge~$0$ (it is a
dimensionful constant, not a field).  Therefore:

First term:
\[
  R[B\,M\,\widetilde{B}]
  = \frac{\Nc}{\Nc+1} + \frac{2}{\Nc+1} + \frac{\Nc}{\Nc+1}
  = \frac{2\Nc + 2}{\Nc + 1} = 2\,.
\]

Second term:
\[
  R[\det M] = (\Nc + 1) \times R[M]
  = (\Nc + 1) \times \frac{2}{\Nc + 1} = 2\,.
\]

Both terms have $R$-charge $2$, and $R[\Lambda^{2\Nc-1}] = 0$, so:
\[
  R[W_{\text{conf}}] = 2\,.
  \quad\checkmark
\]

\bigskip

\textbf{(c) s-Confinement as the boundary of duality.}

The completed table:

\renewcommand{\arraystretch}{1.3}
\begin{center}
\begin{tabular}{@{}l l l@{}}
\toprule
\textbf{Gauge group} & \textbf{s-confining at} &
  \textbf{Dual $\to$ trivial} \\
\midrule
$\SU{\Nc}$ & $\Nf = \Nc + 1$ &
  $\SU{\Nf - \Nc} = \SU{1}$ \\
$\Sp{\Nc}$ & $\Nf = \Nc + 2$ &
  $\Sp{\Nf - \Nc - 2} = \Sp{0}$ \\
$\SO{\Nc}$ & $\Nf = \Nc - 4$ &
  $\SO{\Nf - \Nc + 4} = \SO{0}$ \\
\bottomrule
\end{tabular}
\end{center}

(Recall that for $\Sp{\Nc}$, the matter is $2\Nf$ fundamentals.)

\textbf{Self-dual vs.\ s-confining for $\Sp{1}$.}

With $\Sp{1} = \SU{2}$, the dual is $\Sp{\Nf - 1 - 2} = \Sp{\Nf - 3}$.

\begin{itemize}
  \item $\Nf = 4$ ($2\Nf = 8$ fundamentals):
    $\Sp{4 - 1 - 2} = \Sp{1}$.  This gives back $\Sp{1}$: the theory
    is \textbf{self-dual}, not s-confining.

  \item $\Nf = 3$ ($2\Nf = 6$ fundamentals):
    $\Sp{3 - 1 - 2} = \Sp{0}$ (trivial).  This is the
    \textbf{s-confining} point for $\Sp{1}$.
    \quad$\checkmark$
\end{itemize}

The pattern is general: the self-dual point sits \emph{inside} the
conformal window, while the s-confining point sits at its boundary
(below the conformal window), where the dual gauge group becomes
trivial.

\end{proof}


\bigskip

\begin{center}
\rule{0.6\textwidth}{0.4pt}

\textit{End of Exercise Set 7.  These exercises demonstrate the
breadth of the $\mathcal{N}=1$ SUSY duality landscape: anomaly
matching extends consistently to $\SO{\Nc}$ gauge theories
(Problem~1), the Kutasov--Schwimmer duality provides a family of
dualities parametrised by $k$ that reduces to Seiberg duality at
$k = 1$ (Problem~2), and s-confining theories sit at the boundary
where the magnetic gauge group becomes trivial (Problem~3).  All
three results rest on the same foundation: 't~Hooft anomaly matching
and holomorphy of the superpotential.}
\end{center}


%% ========================================================================
%%  BIBLIOGRAPHY
%% ========================================================================
\begin{thebibliography}{99}

\bibitem{IS9506101}
  K.~Intriligator and N.~Seiberg,
  ``Duality, monopoles, dyons, confinement and oblique confinement
  in supersymmetric $\SO{\Nc}$ gauge theories,''
  \textit{Nucl.\ Phys.} \textbf{B444} (1995) 125
  [\texttt{hep-th/9506101}].

\bibitem{IS1995}
  K.~Intriligator and N.~Seiberg,
  ``Lectures on supersymmetric gauge theories and electric--magnetic
  duality,''
  \textit{Nucl.\ Phys.\ Proc.\ Suppl.} \textbf{45BC} (1996) 1
  [\texttt{hep-th/9509066}].

\bibitem{KS9505004}
  D.~Kutasov and A.~Schwimmer,
  ``On duality in supersymmetric Yang--Mills theory,''
  \textit{Phys.\ Lett.} \textbf{B354} (1995) 315
  [\texttt{hep-th/9505004}].

\bibitem{KSS9507040}
  D.~Kutasov, A.~Schwimmer, and N.~Seiberg,
  ``Chiral rings, singularity theory and electric--magnetic duality,''
  \textit{Nucl.\ Phys.} \textbf{B459} (1996) 455
  [\texttt{hep-th/9510222}].

\bibitem{CSS9612207}
  C.~Cs\'aki, M.~Schmaltz, and W.~Skiba,
  ``A systematic approach to confinement in $\mathcal{N}=1$
  supersymmetric gauge theories,''
  \textit{Phys.\ Rev.\ Lett.} \textbf{78} (1997) 799
  [\texttt{hep-th/9612207}].

\bibitem{Seiberg1995}
  N.~Seiberg,
  ``Electric--magnetic duality in supersymmetric non-Abelian gauge
  theories,''
  \textit{Nucl.\ Phys.} \textbf{B435} (1995) 129
  [\texttt{hep-th/9411149}].

\end{thebibliography}

\end{document}
